\section{Homology through successive layers}\label{sec:homology}
The kernel calculations compare neighbourhoods in the first extension.
We now follow those groups through every layer of the general construction,
with any number of auxiliary types. As in Section~\ref{sec:construction},
write $\mathcal{T}=\{0,\ldots,d-1\}$ for the type set. The coefficients and augmented chain
conventions are those of Section~\ref{sec:kernels}. The degree shift will
supply the criterion used in Theorem~\ref{ker:thm:equivalence}.

\subsection{A degree shift in every non-central neighbourhood}

For complexes on disjoint vertex sets, \(K*J\) consists of unions
of one face from each factor. If \(S^0\) consists of two vertices
without an edge, \(K*S^0\) is the union of two cones on \(K\).
Lemma~\ref{hom:lem:union} gives
\begin{equation}\label{hom:eq:suspension}
 \Ht_t(K*S^0)\cong\Ht_{t-1}(K).
\end{equation}
This includes \(K=\{\varnothing\}\), for which the join is \(S^0\).
The clique complex of the complete join of two graphs is the join
of their clique complexes, since a clique is a union of a clique
in each graph.

\begin{theorem}\label{hom:thm:shift}
For \(m\geq1\) and every integer \(t\),
\begin{align}
 \Ht_t\bigl(\Cl(L_{G_m}(a))\bigr)
 &\cong\Ht_{t-m+1}\bigl(\Cl(L_B(a))\bigr),
 &&a\in V(A),\label{hom:eq:fixed}\\
 \Ht_t\bigl(\Cl(L_{G_m}(x_{j,i}))\bigr)
 &\cong\Ht_{t-m+1}\bigl(\Cl(L_B(x_j))\bigr),
 &&j\in \mathcal{T},\quad 0\leq i<m.\label{hom:eq:auxiliary}
\end{align}
\end{theorem}
\begin{proof}
First fix \(a\in V(A)\) and put \(\Lambda_m=L_{G_m}(a)\).
Its vertices are \(N_A(a)\), all centres, and all copies of the
types \(j\) for which \(a\in P_j\). For \(m\geq2\), delete a centre
\(c_i\) from \(\Lambda_m\). The folds \(x_{j,h}\to x_{j,i}\), \(h\ne i\),
are still valid: restrict the neighbourhood inclusions in the proof of
Lemma~\ref{lem:remainder} to \(N_A(a)\) and the types present.
After these folds, the remaining graph is a cone with a remaining
centre as apex.
If there are no such auxiliary types, it is already a cone.
Lemmas~\ref{hom:lem:cone} and \ref{hom:lem:fold} imply that
\(\Cl(\Lambda_m-c_i)\) is acyclic.

The neighbourhood of \(c_i\) inside \(\Lambda_m\) is isomorphic to
\(\Lambda_{m-1}\), using the same layer reindexing restricted to the
vertices under consideration. Hence the decomposition
\[
 \Cl(\Lambda_m)=\Cl(\Lambda_m-c_i)
       \cup\bigl(c_i*\Cl(L_{\Lambda_m}(c_i))\bigr)
\]
has acyclic pieces and intersection isomorphic to \(\Cl(\Lambda_{m-1})\).
Lemma~\ref{hom:lem:union} shifts homology by one degree. Iteration
proves \eqref{hom:eq:fixed}; at \(m=1\) it is the identity.
The set \(N_A(a)\) may be empty, but \(\Lambda_1\) still contains \(p\).

Now fix an auxiliary vertex \(x_{j,i}\). Its neighbourhood has
the vertices of \(P_j\), all copies of types in \(N_R(j)\),
the vertices \(x_{j,h}\) for \(h\ne i\), and the centres \(c_h\)
for \(h\ne i\). For each \(u\in N_R(j)\), fold
\(x_{u,h}\) onto \(x_{u,i}\) for \(h\ne i\). The source and
target have the same adjacencies to \(P_j\) and auxiliary types.
The target sees all centres in this neighbourhood, whereas
the source may omit \(c_h\). These folds therefore preserve
homology, and their targets remain throughout.

What remains consists of a copy of \(L_B(x_j)\) on
\[
 P_j\ \sqcup\ \{x_{u,i}:u\in N_R(j)\}
\]
and the \(m-1\) pairs \(\{c_h,x_{j,h}\}\), \(h\ne i\).
Each pair has no internal edge. Distinct pairs are completely
joined, and every pair is completely joined to the displayed copy
of \(L_B(x_j)\). The remaining clique complex is consequently
\[
 \Cl(L_B(x_j))*(S^0)^{*(m-1)}.
\]
Apply Lemma~\ref{hom:lem:fold} and then \eqref{hom:eq:suspension}
\(m-1\) times. If \(L_B(x_j)\) is empty, both \(P_j\) and
\(N_R(j)\) are empty. There are then no folds, the first factor is
\(\{\varnothing\}\), and the same suspension identity applies.
For \(m=1\) there are no pairs, so this case too is the identity.
\end{proof}

\subsection{The homological criterion}

Non-zero groups in the first layer remain non-zero after every shift. Together with a reduction of the remainder, they give the following criterion for strictness.

\begin{theorem}\label{hom:thm:criterion}
Let \(s\geq0\). Suppose that
\begin{enumerate}
\item \(A\in\D_s\setminus\D_{s-1}\);
\item \(B-p\in\D_{s+1}\);
\item for every \(v\in V(B)\setminus\{p\}\), some integer \(t\)
satisfies \(\Ht_t(\Cl(L_B(v)))\ne0\).
\end{enumerate}
Then for every \(m\geq0\),
\[
 G_m\in\D_{s+m}\setminus\D_{s+m-1},
 \qquad |V(G_m)|=|V(A)|+(d+1)m.
\]
\end{theorem}
\begin{proof}
Theorem~\ref{hom:thm:shift} propagates non-zero reduced homology to
every non-central neighbourhood in every \(G_m\), \(m\geq1\).
By Lemma~\ref{hom:lem:acyclic}, none of these neighbourhoods belongs
to any class \(\D_j\).

The case \(m=0\) is the first hypothesis. For \(m=1\), the centre
has neighbourhood \(A\in\D_s\) and its deletion leaves
\(B-p\in\D_{s+1}\). Equation~\eqref{eq:first} yields
\(B\in\D_{s+1}\). A first deletion at level \(s\) is impossible:
the centre's neighbourhood is not in \(\D_{s-1}\), and neither is
any non-central neighbourhood. Since \(B\) has more than one
vertex, \eqref{eq:first} gives \(B\notin\D_s\).
This includes \(s=0\), when \(\D_{s-1}\) is the class of cones.

For \(m\geq2\), suppose \(G_{m-1}\) is strictly
\((s+m-1)\)-dismantlable. Each centre has neighbourhood isomorphic
to \(G_{m-1}\); deleting it leaves a graph in
\(\D_0\subseteq\D_{s+m}\) by Lemmas~\ref{lem:centre} and~\ref{lem:remainder}.
Thus \(G_m\in\D_{s+m}\). A first deletion at level \(s+m-1\)
would require a neighbourhood in \(\D_{s+m-2}\).
Centres fail this condition by the induction hypothesis, and
non-central vertices fail it by homology. This proves strictness.
The order follows from \eqref{eq:general-order}.
\end{proof}

\subsection{Comparison with Euler characteristic}

The obstruction \(q\) of Section~\ref{sec:construction} is the alternating
sum of reduced Betti numbers, with a contribution from degree \(-1\).
Rank--nullity in the augmented chain complex gives
\begin{equation}\label{hom:eq:euler}
 q(H)=\sum_{r\geq0}(-1)^r
       \dim_{\F}\Ht_{r-1}(\Cl(H)).
\end{equation}
Indeed, if \(f_t=\dim\Ct_t\), \(b_t=\dim\Ht_t\), and
\(\rho_t=\operatorname{rank}\partial_t\), then
\(f_t=b_t+\rho_t+\rho_{t+1}\). In the finite alternating sum
the ranks cancel, including the term of degree \(-1\).
In particular, \(q(H)\ne0\) implies non-trivial reduced homology.

The converse need not hold. The graph \(C_4\sqcup K_1\) has five
vertices, four edges and no triangles, so \(q=1-5+4=0\).
Its two components give \(\dim\Ht_0=1\). The boundary matrix of
the four cycle edges has rank three, and there are no
two-dimensional faces, giving \(\dim\Ht_1=1\).
These contributions cancel in \eqref{hom:eq:euler}.
This example separates the two obstructions on a graph; it does
not by itself supply initial data satisfying all three hypotheses
of Theorem~\ref{hom:thm:criterion}.

\Needspace{10\baselineskip}
The Euler criterion in Theorem~\ref{thm:template} uses
\(q(L_B(v))\ne0\) for every non-central \(v\), together with
\(A\in\D_s\setminus\D_{s-1}\) and the stronger remainder condition
\(B-p\in\D_s\). Equation~\eqref{hom:eq:euler} and
Lemma~\ref{lem:inclusions} imply the hypotheses of Theorem~\ref{hom:thm:criterion}.
Allowing the remainder in \(\D_{s+1}\) is a separate change from
replacing Euler characteristic by homology: the first-deletion
argument at \(m=1\) also permits this weaker remainder in the
Euler proof. Both criteria retain a distinct requirement on the
remainder; non-trivial homology of the other neighbourhoods alone
does not ensure the positive reduction.


\runcheck{homology}
The calculation checks the empty graph, a cone, $C_4\sqcup K_1$, and an
auxiliary with an empty first-layer neighbourhood. In the last example,
after $m$ layers its only non-zero reduced Betti number is $\widetilde\beta_{m-2}=1$,
for the tested values $m=1,2,3$. The full command
\texttt{python reproduce.py --extended} additionally checks all labelled
seeds on one, two or three vertices, all profiles for one or two types,
and all type graphs, at $m=0,1,2,3$. Homology is computed for the subcollection
with seed order at most two. These are finite checks of the identities;
Theorem~\ref{hom:thm:shift} proves them for arbitrary seeds and all layers.
